% 02-Background.tex. All sections \framedecl{-}. No claim blocks here (all
% twelve are in chapter 4). One \needsaudit, in section 2.3.

\chapter{Background}
\label{sec:litreview}

\noindent
This chapter is not a survey, but the shortest route to what this field measures,
what those measurements assume, and where they give way.

% ===========================================================================
\section{Transformers on single-cell data}
\label{sec:lit-foundation}
\framedecl{-}

\begin{question}
A recipe built for text is now used to write cells. What made it portable, and
what does adopting it commit you to? It asks almost nothing of language, and
everything it does ask falls on how the data are written down.
\end{question}

\noindent
The transformer recipe requires only sequences of discrete tokens, some
predictable from others. Nothing there mentions words, which invited the transfer
into biology. scGPT~\cite{scgpt} and Geneformer~\cite{geneformer} pretrain on
tens of millions of cells written as gene-token sequences, toward some version of
the \emph{virtual cell}~\cite{bunne}, whose simulated responses could substitute
for an experiment.

That sets the bar for everything measured later, but it is not where the models
differ. A cell is a vector over some twenty thousand genes; a transformer
consumes discrete tokens. What must be discarded in between sets a bound on what the model can learn.

% ===========================================================================
\section{Cell2Sentence and C2S-Scale}
\label{sec:lit-c2s}
\framedecl{-}

\begin{question}
If a cell has to be written as tokens, something must go. Which choice does the
encoding make, and what argument says the loss is safe? Rank is kept, magnitude
is thrown away, and the defence is a reconstruction argument.
\end{question}

\noindent
Cell2Sentence~\cite{c2s} takes the most literal route available. Rank a cell's
genes by expression, write the symbols out in descending order, and the cell is a
string of ordinary text, a \emph{cell sentence}. Rank is the whole of the
representation, which symbols appear and in what order. Magnitude is not written
down at all.

% The ten most abundant panel genes by mean measured expression over the 606
% consensus truth profiles (nir_consensus_profiles.npz), in rank order, so every
% gene shown occurs in this thesis's data. It previously showed a
% whole-transcriptome sentence (MALAT1 ACTB TMSB4X B2M EEF1A1 RPL13A FTL TPT1
% RPS27 GAPDH), nine of whose ten genes are not in the 946-gene panel.
\begin{quote}
\ttfamily\small
GAPDH RPS6 KTN1 EXT1 PTK2 MACF1 HSPD1 ALDOA GNAS CD44 \dots
\end{quote}

The same literalness extends to everything the model is given. Because a cell
sentence is ordinary text, the prompt that surrounds it is ordinary text too
(\cref{fig:prompt}): one English instruction naming the cell line, the drug,
the dose and the mechanism annotation, then the untreated cell written as a
sentence, then an open-ended cue that the model completes with a sentence of its
own. Nothing about the cell is handed to the model as a vector, and no part of
the architecture is reserved for it. That is what makes an off-the-shelf
language model usable here without modification, and it is the property the
encoding buys: the model can bring to bear whatever it learned about drug names,
mechanisms and gene symbols in natural-language pretraining, and it reasons
about the cell's sequencing in the same medium it reasons about anything else,
as a sequence of words. Which string it is trained to complete the prompt with
is the one thing that varies across this thesis, and \cref{fig:prompt} shows the
two it is asked for.

\begin{figure}[htbp]
% MARGINALIA: fits the text column; caption in the margin, as fig:two-questions.
\begin{lrbox}{\tvfb}%
% figs/prompt.tex -- the prompt as the model receives it, and the two strings it
% is trained to continue it with. In-column figure, caption in the margin
% (\sidecap), same idiom as figs/two-questions.tex. The instruction line is the
% literal template of shared/evaluate_c2s_tahoe.py and
% endcell/preprocess/tahoe_c2s_preprocess_endcell_v2.py::format_prompt; cell
% line, drug, dose and mechanism are real atlas values (CVCL_0480 = PANC-1,
% Docetaxel, 0.05 uM, moa-fine "Microtubule inhibitor"). The gene lists are the
% panel sentence quoted in sec:lit-c2s and the illustrative residual of
% fig:frame, truncated; they carry no measurement.
\begin{tikzpicture}[
  x=1cm, y=1cm, font=\small,
  lbl/.style={font=\scriptsize\bfseries\sffamily, text=accent, anchor=south west},
  note/.style={font=\scriptsize, text=quietink, anchor=north west},
  box/.style={draw=rulegrey, rounded corners=2pt, inner sep=6pt,
              fill=panelbg, align=left, anchor=north west},
  ar/.style={-{Stealth[length=4.5pt]}, quietink, line width=0.6pt},
]
\newcommand{\mono}[1]{{\ttfamily\scriptsize #1}}
\newcommand{\spanhl}[1]{\colorbox{rulegrey!45}{#1}}

% --- the prompt ------------------------------------------------------------
\node[lbl] at (0,0) {PROMPT};
\node[box, text width=11.6cm] (prompt) at (0,-0.08) {%
  \mono{Predict the response of \spanhl{PANC-1} to \spanhl{Docetaxel} at
  \spanhl{0.05 uM}. Mechanism: \spanhl{Microtubule inhibitor}.}\\[2pt]
  \mono{Control cell: GAPDH RPS6 KTN1 EXT1 PTK2 MACF1 HSPD1 ALDOA GNAS CD44
  \ldots\ [END\_CELL]}\\[6pt]
  \mono{Response cell:}};
\node[note, text width=11.6cm] (note1) at ($(prompt.south west)+(0,-0.06)$)
  {one instruction line in English, the untreated cell as a sentence, and an
   open-ended cue the model completes};

% --- the arrow -------------------------------------------------------------
\coordinate (y0) at ($(note1.south west)+(0,-0.10)$);
\draw[ar] ($(y0)+(5.95,0)$) -- ($(y0)+(5.95,-0.72)$);
\node[font=\scriptsize\itshape, text=ink, anchor=west, text width=5.7cm]
  at ($(y0)+(6.25,-0.36)$)
  {the model continues the text; the target decides with what};

% --- the two targets -------------------------------------------------------
\node[lbl] at ($(y0)+(0,-1.30)$) {STANDARD TARGET \textcolor{quietink}{(\Eframe)}};
\node[box, text width=5.4cm] (ta) at ($(y0)+(0,-1.38)$) {%
  \mono{GAPDH RPS6 HSPD1 CDKN1A ALDOA GNAS \ldots\ [END\_CELL]}};
\node[note, text width=5.5cm] at ($(ta.south west)+(0,-0.06)$)
  {the treated cell's own sentence: ranked genes, most abundant first};

\node[lbl] at ($(y0)+(6.2,-1.30)$) {RESIDUAL TARGET \textcolor{quietink}{(\Rframe)}};
\node[box, text width=5.4cm] (tb) at ($(y0)+(6.2,-1.38)$) {%
  \mono{CDKN1A DDB2 MDM2 \ldots\ [DOWN] HMOX1 TXNIP \ldots\ [END\_CELL]}};
\node[note, text width=5.5cm] at ($(tb.south west)+(0,-0.06)$)
  {the drug-specific residual of \cref{sec:apparatus-vocab}: up-regulated
   genes, a separator, down-regulated genes};
\end{tikzpicture}
%
\end{lrbox}
\sidecap[The prompt and its two targets]{figure}{\captiontitle{The model
sees plain text and returns plain text.} The prompt is the literal template
every arm uses: an instruction line whose four shaded spans carry the cell line,
the drug, the dose and the mechanism annotation, the control cell as a sentence
closed by \ENDCELL{}, and the cue \texttt{Response cell:}. The four values are
real atlas entries; the gene lists are illustrative and carry no measurement.
Below, the two strings the model is trained to continue with: the treated cell's
own sentence in the expression frame, and the signed drug-specific residual
of \cref{sec:apparatus-vocab} in the residual frame, whose two blocks are split
by \DOWN{}. Both frames share the prompt to the byte; only the target
differs.\label{fig:prompt}}
\end{figure}


\paragraph{What the defence covers} C2S defends the discarded magnitude by
reconstruction. Expression can be recovered from rank through a fitted
log-linear relationship $e_i = a_d\log(r_i) + b_d$, achieving $R^2 = 0.815$
against the original values on peripheral blood data. Rank, they argue,
therefore encodes expression reversibly enough for the loss not to matter
downstream.

The criterion is absolute, not discriminative. Reconstruction fidelity asks how
close the recovered profile is to the true one, which is not the same as
preserving what distinguishes one condition from another. That distinction runs
through the whole thesis. The map is not only C2S's defence of the discarded
magnitude but also part of this thesis's own measurement path: in the expression
frame a prediction and every truth it is scored against are decoded through the
same fitted relationship before any distance between them is taken
(\cref{sec:res-metrics}).

\paragraph{Three properties that matter later}
C2S-Scale builds on decoder-only transformers of the GPT family, the checkpoint
used here on Pythia~\cite{pythia}. Token embeddings form a \emph{residual
stream}, one vector per position that every layer reads from and writes back into
additively, so a representation can be read at depth and
modified.\seealso{\cref{sec:res-mechanism} performs the intervention this design
permits} Training weights every token position equally, so a target largely
shared between conditions gives any one condition only a small share of the
gradient. And conditioning runs entirely through attention across a causal mask.

\paragraph{The family and the checkpoint}
C2S-Scale~\cite{c2sscale} spans backbones from Pythia at 410 million and 1
billion parameters up to Gemma-2 at 27 billion, pretrained on over 50 million
human and mouse cells drawn from CELLxGENE and the Human Cell Atlas. None of
that pretraining is a perturbation task: the objectives are cell sentence
generation, cell type and tissue prediction, and text tasks over the
accompanying metadata. Perturbation prediction enters only afterwards, as a
separate fine-tuning stage, and it is evaluated on two benchmarks that ask
different questions.

The first is single-cell. Immune cells are exposed to cytokines; the model
receives an untreated cell's sentence together with a cell type, a cytokine and
an exposure duration, and must generate the treated cell's sentence while
recovering those three labels. Held-out conditions are graded into difficulty
tiers by how far they sit from what training contained.

The second is bulk, and it is a transfer test across cell lines rather than
across compounds. L1000 reads a fixed panel of 978 landmark genes from a
treated well, one profile per compound and cell line pair. The profiles are
quantile-normalised, and for each compound half of the cell lines are held out.
The compound itself is therefore present in training, applied to other lines,
and the question is whether the model can carry its effect onto a line it has
never seen treated with it.

Two features of that design bear on the comparison later. Pairing is by group
and not by cell: in the bulk arm treated and untreated profiles are matched on
cell line name alone, which is unavoidable once sequencing destroys the cell it
reads. And the single-cell results are scored as distributions, generated
population against observed and assessed collectively rather than condition by
condition, which asks whether the model produces a plausible population and
never whether it tells one condition apart from another.

The checkpoint fine-tuned below, \texttt{C2S-Scale-Pythia-1b-pt}, is the
pretrained model rather than either of those fine-tunes, and no perturbation
fine-tune of the family has been released. The starting point here is the
authors' own, with none of their perturbation training already in it.
% ===========================================================================
\section{Latent-space perturbation models}
\label{sec:lit-chem}
\framedecl{-}

\begin{question}
The obvious way to predict a treated cell is not this one. What do the standard
methods do differently? In every one except the model this thesis fine-tunes, the
drug arrives through a channel the architecture was built to receive it on.
\end{question}

\noindent
The direct approach predicts a treated cell as a displacement in a learned
latent space. scGen~\cite{scgen} estimates the vector between control and
perturbed populations in a variational autoencoder's latent space and adds it to
a context where the perturbation was never applied. CPA~\cite{cpa} makes that
arithmetic learnable, composing a basal cell state with additive embeddings for
drug, dose and covariates; chemCPA swaps its one-hot drug embedding for a
molecular-structure encoder, letting an unseen drug be embedded at
all.\footnote{Hetzel et al., \emph{Predicting Cellular Responses to Novel Drug
Perturbations at a Single-Cell Resolution}, NeurIPS 2022.} GEARS~\cite{gears}
reaches unseen genes through a gene--gene knowledge graph. CellOT instead fits a
neural optimal-transport map from control to treated per perturbation.\footnote{
Bunne et al., \emph{Learning single-cell perturbation responses using neural
optimal transport}, Nature Methods 2023; Bunne, Krause and Cuturi,
\emph{Supervised Training of Conditional Monge Maps}, NeurIPS 2022.}

\paragraph{Two properties they share} The prediction is a vector in a continuous
space, emitted in one pass by a decoder trained on a loss written directly on
that vector. And the drug enters on a channel built to receive it, an embedding
in CPA, chemCPA and GEARS, a latent offset in scGen, a fitted map in CellOT.

A cell-sentence model does neither. It samples a token sequence left to right,
carries the drug as a few words of prompt, and weights every token position
equally in its loss. In the standard cell-sentence target, $83\%$ of the
top-$200$ gene positions are identical between any two drugs in the same context
(\cref{sec:methods-residual}); a gene spans several sub-tokens, so that count is
not a count over loss positions. The string to be emitted is nearly the same string
whichever drug was named, and a latent-transition model has no target with that
property.

Whether that property is the binding constraint is not settled here.
\cref{sec:res-encoding} reports what re-encoding the target does, and stops
short of establishing that tokenisation is what binds.

\paragraph{Scored differently} The latent-transition literature is evaluated
condition by condition against the observed population: distributional distances,
$R^2$ on pseudobulk, DE-gene overlap. \Cref{fig:two-questions} draws where that
diverges from a discriminative question.

\begin{figure}[htbp]
% MARGINALIA (06-figures.md): fits the text column; caption in the margin.
\begin{lrbox}{\tvfb}%
% Body moved out to figs/two-questions.tex. One change from the inline v1
% version: the RIGHT panel's $T$ label is displaced 2pt perpendicular to the
% P--Ta segment so no rule runs under the glyph. Points, caption and label are
% untouched; \cloudpts is still one macro instantiated in both scopes, which is
% what makes "Point positions are identical in both panels" structurally true.
% ===========================================================================
% fig:two-questions -- condition-referenced vs discriminative.
% Figure BODY only. \input this inside a figure environment; the caption and
% \label stay in Sections/02-Background.tex so the surrounding prose owns them.
% Requires: tikz + calc, and the preamble's colours (accent, ink, quietink,
% rulegrey, panelbg). No data. Schematic; no axis carries a value.
%
% Drawn width 402.9pt against a 404.0pt measure -- it is drawn at size and must
% NOT be rescaled by \resizebox or a width= key.
%
% CHANGED FROM v1 -- exactly one thing:
%   In the RIGHT panel the $T$ node label sat on the accent segment P--Ta, so
%   the 0.9pt rule ran under the glyph. The label is displaced 2pt PERPENDICULAR
%   to that segment, away from it. Nothing else moved.
%   P = (3.03,2.69), Ta = (2.93,3.13);  d = Ta-P = (-0.10,0.44), |d| = 0.4512205
%   unit u  = (-0.2216233, 0.9751425)
%   normal n = ( 0.9751425, 0.2216233)   -- the side the label already sits on
%              (old offset (0.10,0.03) has n-component +0.10416 > 0), so +n is
%              "away from the segment".
%   2pt = 0.07029196 cm  ->  delta = (0.06854, 0.01558)
%   old offset (0.10000, 0.03000)  ->  new offset (0.16854, 0.04558)
%   The LEFT panel keeps the original (0.10,0.03): no rule runs under it there.
%
% DO NOT touch \cloudpts or split it per panel. Both scopes instantiate the one
% macro, which is what makes the caption's claim "Point positions are identical
% in both panels" structurally true rather than merely checked.
% ===========================================================================
\begin{tikzpicture}[font=\small, x=1cm, y=1cm]
  % Five points defined ONCE, instantiated in both panels, so they cannot drift.
  \newcommand{\cloudpts}{%
    \coordinate (Ta) at (2.93,3.13);
    \coordinate (Tb) at (3.90,2.31);
    \coordinate (Tc) at (2.04,2.08);
    \coordinate (Td) at (3.26,3.74);
    \coordinate (P)  at (3.03,2.69);}

  \begin{scope}
    \node[anchor=south, font=\scriptsize\bfseries\sffamily, text=accent]
      at (3.0,5.42) {CONDITION-REFERENCED};
    \node[anchor=north, font=\scriptsize\itshape, text=ink, align=center,
          text width=5.4cm]
      at (3.0,5.30) {does the prediction resemble\\ what its own condition
                     produced?};
    \fill[panelbg] (3.0,2.75) ellipse (1.42 and 1.20);
    \cloudpts
    \foreach \n in {Tb,Tc,Td}{\draw[rulegrey, line width=0.6pt]
      (\n) circle (2.0pt);}
    \draw[accent, line width=0.9pt] (P) -- (Ta);
    \fill[ink] (Ta) circle (2.0pt);
    \fill[accent] (P) circle (2.2pt);
    \node[anchor=south west, font=\scriptsize, text=ink]
      at ($(Ta)+(0.10,0.03)$) {$T$};
    \node[anchor=north east, font=\scriptsize, text=accent]
      at ($(P)+(-0.06,-0.04)$) {$P$};
    \node[anchor=north, font=\scriptsize, text=quietink, align=center,
          text width=5.4cm]
      at (3.0,1.14) {consults one truth};
  \end{scope}

  \draw[rulegrey, line width=0.4pt] (6.20,0.82) -- (6.20,5.20);

  \begin{scope}[xshift=6.4cm]
    \node[anchor=south, font=\scriptsize\bfseries\sffamily, text=accent]
      at (3.0,5.42) {DISCRIMINATIVE};
    \node[anchor=north, font=\scriptsize\itshape, text=ink, align=center,
          text width=5.4cm]
      at (3.0,5.30) {does it resemble its own condition\\ more than some other
                     drug's?};
    \fill[panelbg] (3.0,2.75) ellipse (1.42 and 1.20);
    \cloudpts
    \foreach \n in {Tb,Tc,Td}{%
      \draw[quietink, densely dashed, line width=0.45pt] (P) -- (\n);}
    \draw[accent, line width=0.9pt] (P) -- (Ta);
    \foreach \n in {Tb,Tc,Td}{\fill[ink] (\n) circle (2.0pt);}
    \fill[ink] (Ta) circle (2.0pt);
    \fill[accent] (P) circle (2.2pt);
    % 2pt perpendicular to P--Ta, away from it. See header for the arithmetic.
    \node[anchor=south west, font=\scriptsize, text=ink]
      at ($(Ta)+(0.16854,0.04558)$) {$T$};
    \node[anchor=north east, font=\scriptsize, text=accent]
      at ($(P)+(-0.06,-0.04)$) {$P$};
    \node[anchor=north, font=\scriptsize, text=quietink, align=center,
          text width=5.4cm]
      at (3.0,1.14) {consults the whole candidate set};
  \end{scope}
\end{tikzpicture}
%
\end{lrbox}
\sidecap[Two questions about one prediction]{figure}{\captiontitle{The same prediction, the
same truths, two different questions.} $P$ is a prediction for one condition, $T$
is that condition's own truth, and the other three points are rival candidates,
all inside the pale region when the part of the shift shared across drugs
dominates. A
condition-referenced metric consults $T$ alone, so the rivals are hollow on the
left; a discriminative metric grades where $T$ ranks among them. Point positions
are identical in both panels. No axis carries a value.\label{fig:two-questions}}

\end{figure}

\paragraph{Where the field has consolidated}
scPerturb~\cite{scperturb} harmonises the datasets and standardises distances
between populations, and PerturBench~\cite{perturbench} supplies the
exception this thesis relies on, a rank-based score grading a prediction against
a set of candidate conditions and not one truth. None of these is run as a
baseline here; the results chapter concedes the missing head-to-head as a real
gap.

\begin{scopelimit}[carried into \cref{sec:res-lookup} and \cref{sec:limitations}]
What is reported below is not evidence about CPA, chemCPA or CellOT, none of
which was run here. The architectures named in this section supply the
\emph{motivation} for the baseline built in \cref{sec:res-lookup} --- they tell
us what a drug-specific channel typically carries --- and they are not
competitors that have been scored and beaten.
\end{scopelimit}

\paragraph{What is offered instead} Something narrower, measured and not argued.
In scGen and in CPA the drug-specific content carried into a new context is a
single per-drug vector, dose-scaled in CPA's case. \cref{sec:res-lookup}
measures what such a vector is worth by looking the drug up: the mean residual of
the same drug in other cell lines, fitted from training conditions
only.\framecue{\Rframe}{the scores below are residual-frame; see
\cref{sec:methods-residual}}

On the conditions where both are defined it scores $0.913$(NIR). On the held-out wells, where the lookup cannot
borrow the target's own well, it scores $0.890$(NIR). Chance is $0.50$ under the
exchangeability condition audited in \cref{sec:res-metrics}, so both comparisons
sit far from chance, and what the section turns
on is the ordering, not the size of either support.

That is not a score for CPA or chemCPA, whose decoders make the emitted shift
context-dependent even when the drug channel is one vector. But it does say which
question is informative. Not which architecture wins. Whether anything improves
on a per-drug constant.\seealso{\cref{sec:res-lookup} builds the ladder these
scores sit on}


% ===========================================================================
\section{Baselines and metric calibration}
\label{sec:lit-dispute}
\framedecl{-}

\begin{question}
Does any of these deep learning models work? The field has answered both ways, and what decides which
answer to believe is not another model but an audit of the metrics. Both camps
agree the standard metrics are defective; they disagree about what a repaired
measurement would reveal.
\end{question}

\paragraph{One side, trivial predictors are not beaten}
Ahlmann-Eltze, Huber and Anders~\cite{ahlmanneltze} find that deep-learning
prediction of genetic perturbation effects does not outperform linear baselines.
Csendes et al.~\cite{csendes} reach the same conclusion for foundation cell
models, Bendidi et al.~\cite{bendidi} report that a single PCA still rules them
all, and Kernfeld et al.~\cite{kernfeld} concur across expression-forecasting
methods. Nor is this confined to transcriptomic readouts. Hauptmann and
Kramer~\cite{hauptmann} hold the protocol fixed and watch the differences between
architectures dissolve; DrEval~\cite{dreval} finds deep models barely beating a
mean drug and cell-line predictor.

\paragraph{The other side, the metrics were at fault} Miller et
al.~\cite{miller} argue that these conclusions ``largely stem from limitations of
benchmarking metrics, not from the models themselves''. They formalise positive
and negative controls in perturbation benchmarking and introduce the
\emph{interpolated duplicate}, blending a technical duplicate with the mean
baseline by per-gene differential-expression
significance.\gloss{interpolated duplicate}{A positive control of known and
tunable quality: a technical duplicate blended with the mean baseline.}

\begin{defn}{Dynamic range fraction (\DRF)}
A quantitative measure of whether a metric can see quality at all. Given
reference points whose quality is known and tunable --- the interpolated
duplicate supplies them --- \DRF{} grades the metric on how much of that known
difference in quality it registers. It is a property of the metric, not of any
model scored with it.
\end{defn}

Miller et al.\ apply it across 14 Perturb-seq datasets and 13 metrics. MSE and
control-referenced Pearson correlation come out poorly calibrated; weighted MSE,
weighted $R^2\Delta$~\cite{mejia} and Normalised Inverse
Rank~\cite{perturbench} are consistently calibrated. On the metrics that survive,
they report, deep models \emph{do} beat mean, control and linear baselines.

\paragraph{A third result, on the same side of the field}
Vi\~nas Torn\'e et al.~\cite{systema} agree with the first camp about what is
found and with the second about how to look, which is what makes their result
more than another vote. Across ten genetic-perturbation datasets --- three
technologies, five cell lines, from $81$ to $\num{1813}$ perturbations --- two
trivial baselines match or beat CPA, GEARS and scGPT. The \emph{perturbed mean}
predicts the average of every perturbed cell in the training set, the same
profile for all perturbations, and it wins on Pearson-$\Delta$ across every
dataset they test.

The diagnosis is what matters here. Those baselines win because perturbed and
control cells differ in a direction that is shared across perturbations, so a
prediction reproducing only the shared direction already scores well against a
control-referenced target.

\begin{defn}{Systematic variation}
The consistent transcriptional difference between perturbed and control cells
that is common to perturbations rather than specific to any one of them. Vi\~nas
Torn\'e et al.\ measure it as the cosine between a perturbation's own shift from
the control centroid and the average perturbation effect, averaged over
perturbations. It is a property of a dataset, not of a model.
\end{defn}

In their data it is neither small nor incidental. It ranges from about $0.2$ to
$0.7$ across the ten datasets and traces to concrete causes: panel selection ---
Norman et al.~\cite{norman} chose genes in cell-cycle and erythroid pathways ---
and the perturbation technology itself, with CRISPR-induced chromosomal
instability placing $46\%$ of perturbed RPE1 cells in G1 against $25\%$ of
controls. Across datasets the measured systematic variation tracks reported
performance at $r = 0.91$ for GEARS and $0.95$ for fine-tuned scGPT. To that
extent a benchmark ranking is reporting a property of the data.

Their repair is one change, and it is a change to the measurement rather than to
any model: score against the \emph{perturbed centroid}, the average of the
per-perturbation centroids, instead of against the control centroid. On that
reference the numbers collapse --- scGPT on Norman from Pearson$\Delta$ $0.69$ to
$0.13$, GEARS from $0.65$ to $0.13$ --- and their correlation with systematic
variation goes with them. Repairing the reference did not vindicate the models.
It is again, though, a genetic-perturbation answer.

\paragraph{What is actually in dispute} One reading is that the models were
always capturing real biology and the broken metrics hid it; the other is that
they were never capturing much, and better measurement will show that failure
more clearly.

Either way the question is empirical, and settled only in the setting where it is
tested. Miller et al.\ and Ahlmann-Eltze et al.\ both study \textbf{genetic}
perturbation, a designed near-binary knockout with a named target gene. DrEval
studies \textbf{drug sensitivity} as a scalar potency such as IC$_{50}$. None
takes as its object a full transcriptomic response to a drug, which acts at a
dose and engages several proteins at once, at single-cell resolution, from an
autoregressive generative model. The models of \cref{sec:lit-chem} do work
there. What we are not aware of anyone having done is to run the calibration audit first, letting it decide which metrics are admissible before any model is
scored.

\paragraph{The position taken here}
\cref{sec:res-metrics} ports the \DRF{} framework onto drug perturbation, so the
verdict cannot be dismissed by the objection Miller et al.\ themselves raise. The
framework and its interpolated-duplicate positive control are inherited; not
every choice inside them is. The surviving metric uses a similarity function we
picked, Euclidean distance in expression space and not over rank vectors, and the
two variants disagree. The family graded is our selection of five, including our
own implementation of their weighted $R^2$ variant.

\paragraph{Two separations from Miller and colleagues} Only the first is a metric-level
disagreement, and it begins with weighted $R^2\Delta$, which
they endorse as calibrated but which is negative in both arms under our
implementation. Its sign tracks cell count, so \cref{sec:res-metrics} draws no
directional conclusion from that negative, but this member of their calibrated
set does not transfer intact. Their endorsed $\NIR$ does survive. $\paneltau$ and
\texttt{spearman\_expr}, neither in that set, both invert in this field-wide
audit; that failure does not generalise to $\NIR$ or to rank-based metrics as a
family.

The second separation concerns models. Miller et al.\ report that deep models
beat uninformative baselines once the metric is fixed; on the metric their
analysis endorses, the checkpoint evaluated here does not
(\cref{sec:res-blind}). Both can stand. A repaired metric is a precondition for
seeing skill, not a guarantee that there is skill to see. Which difference
accounts for the split, genetic against chemical perturbation, regression against
generation, or the per-drug baseline, is taken up in \cref{sec:res-metrics}
and~\cref{sec:res-lookup}.

\paragraph{The target this thesis inherits from Systema} 

\Cref{sec:methods-residual} is Vi\~nas Torn\'e et al.'s reference change, carried
into the training target as well as the score. Carrying it there invites an objection better answered before the construction is
used than after. A checkpoint trained to emit residuals is then scored in the
residual frame, which is to say graded in the representation it was optimised
for. This looks like a problem and is not one, because the claim resting on this frame is not a
level but a gap: one checkpoint against its own scrambled prompt, on the same
conditions, in the same frame, with the same generic subtracted and the same
scoring rule applied, differing only in which drug the prompt names. Whatever a
shared frame confers is present in both terms and subtracts out. The metric
supplies that protection a second time, since $\NIR$ ranks a prediction against
the other drugs of its own group and any component common to that comparison set
cancels from the ranking --- the generic among them, and the error in estimating
it, because the generic is fitted on exactly the group the comparison set is
drawn from. A model that had learned what a residual looks like and nothing about
which drug it was named would sit at $0.500$ and not above it. What the shared
frame does buy is precision and not position: the drug-specific part is a larger
share of what is being compared, so the same ranking is read with less noise.
\Cref{sec:res-metrics} fixes the metric before any arm is scored, and
\cref{sec:res-encoding} states the scope the reliability filter imposes on the
population it is read over.


Their perturbation-specific effect is
the perturbation's own centroid measured against the average of the
per-perturbation centroids, fitted over the training perturbations,
\begin{equation}
 \Delta_X \;=\; O(X) - O_{\textrm{pert}},
 \qquad
 O_{\textrm{pert}} \;=\; \frac{1}{|\mathcal{P}_{\textrm{train}}|}
   \sum_{X' \in \mathcal{P}_{\textrm{train}}} O(X') .
\end{equation}
The residual target we use is written as a control-referenced shift with a generic
subtracted, $\resid(d,c) = (\bar{y}_{d,c} - \bar{y}_{\textrm{ctrl},c}) -
\textrm{generic}(c)$, but the control cancels and what is left is the same
object,
\begin{equation}
 \resid(d,c) \;=\; \bar{y}_{d,c} \;-\; \frac{1}{D-1}\sum_{d' \neq d} \bar{y}_{d',c} .
\end{equation}
Two further choices are theirs as well. The reference averages over
\emph{perturbations} and not over cells, so a drug run at five doses does not
carry five times the weight of a drug run at one --- the same distinction they
draw between $O_{\textrm{pert}}$ and the cell-weighted perturbed mean they use as
a baseline. And it is fitted on the training split only, which is what
\cref{sec:methods-residual} enforces as split-before-fit.

\paragraph{Three separations from Vi\~nas Torn\'e and colleagues}
Each is forced by the
setting rather than preferred, and the setting differs in three ways that matter:
many cell lines instead of one, scoring on trained conditions as well as
held-out ones, and a generative model that emits text rather than a vector.

The first is scope. Their reference is one centroid for a dataset, which is
well specified when the dataset is one cell line and one screen; Norman is K562,
Replogle is RPE1 or K562. This atlas is $\num{50}$ cell lines across many plates,
and the shared response is context-dependent: a cytotoxic dose does not leave the
same common signature in every line. A single global reference would subtract the
wrong quantity nearly everywhere and leave cell-line and plate structure
\emph{inside} the residual --- which is disqualifying here in particular, because
$\NIR$ ranks a drug against other drugs \emph{within} a context, so retained
context structure would survive into the comparison the metric makes. The generic
is therefore estimated within a (cell line, plate) group.

The second is that the drug is always excluded from its own reference. Their
$O_{\textrm{pert}}$ runs over training perturbations while their evaluation is on
held-out ones, so the perturbation being scored is excluded already. This thesis
scores trained conditions too (\cref{sec:res-encoding}), and without an explicit
leave-one-drug-out a trained condition would be centred on a mean containing
itself while a held-out condition is not, so the two splits would be scored
against differently defined targets. Since the comparison between them is the
claim (\cref{sec:res-generalisation}), the two must be built the same way. The
bias removed is of order $1/D$, which is larger here than in their datasets
because a plate group holds fewer drugs than their panels hold perturbations.

The third is where the correction is applied, and it is the substantive one.
Systema is an evaluation framework: the reference moves at scoring time and the
models are left alone. That is sufficient when the question is whether reported
scores are inflated. It is not sufficient here, for two reasons.
\Cref{sec:methods-residual} finds that the shared component does not only inflate
the score but dominates the \emph{target}: $83\%$ of the top-$200$ gene positions of a standard target string are identical between any two drugs in the same context,
against $40\%$ under the residual. A model trained on a target that barely
encodes the drug cannot be repaired by rescoring its output. And a cell sentence
is a rank order of gene names, not a vector, so it does not support subtraction
at all; recentring it at scoring time would require decoding back to expression
through a channel that keeps rank and discards magnitude
(\cref{sec:lit-c2s}). The arithmetic has to happen where arithmetic is
still defined, which is on the target before it is tokenised. 

\paragraph{A limitation of theirs, and the filter it motivates} Vi\~nas Torn\'e
et al.\ note that a reference-sensitive metric is poorly suited to weak
perturbations, because the shift has small norm and a few cells can swing its
orientation, and they leave this open. Moving the reference into the target
inherits the problem and sharpens it: a residual too small to estimate stably is
not merely scored noisily, it is \emph{trained on}. \Cref{sec:methods-residual}
answers it with a reliability criterion --- a condition enters only if its
residual, estimated on two disjoint halves of its own cells, reproduces above a
threshold fixed at a matched null. That filter is a measurement-quality gate
inherited from their limitation, not a performance one, and its retention rates
are reported per split so the selection stays visible.

\paragraph{How much systematic variation this atlas carries}
Their statistic is defined on a dataset, so it could in principle be computed
here and read on the same scale as their ten. It was not computed for this
thesis, and no value is quoted for it. What this thesis measures instead is the
share of the treated-minus-control shift's squared norm that the generic
programme carries, $0.390$ across the cache (\cref{sec:res-transfer}), read
within the (cell line, plate) group that the first separation above argues the
reference must be scoped to. That number matters for the argument rather than
decorating it: the case for re-encoding the target rests on the shared component
being large in \emph{this} data, and that is a measurement, not an inheritance
from their datasets.

% ===========================================================================
\section{DrEval: pseudoreplication, bias and splits}
\label{sec:lit-dreval}
\framedecl{-}

\begin{question}
A benchmark can report a number that is arithmetically correct and still mislead
about everything a reader cares about. Which failure modes are catalogued? Four
of DrEval's six shape the apparatus directly.
\end{question}

\noindent
DrEval~\cite{dreval} identifies six obstacles to clinical translation:
non-reproducible models, data leakage producing weak generalisation,
pseudoreplication, biased evaluation metrics, missing ablation studies, and
inconsistent response data.

\paragraph{Pseudoreplication} In Hurlbert's original sense~\cite{hurlbert}, this
is treating repeated measurements of one unit as if they were independent
replicates. GDSC2's $395{,}025$ measurements descend from only $969$ cell lines
and $287$ drugs, so an interval computed over measurements is badly
overconfident.

Our data has the same shape, hundreds of thousands of cells resolving to a few
thousand conditions across $50$ cell lines, and no interval here is computed over
cells. Three estimators stand in their place. An interval on a mean is clustered
on the two units that repeat and are crossed, the cell line and the treatment
well. A statistic over \emph{pairs} of conditions uses a dyadic estimator, since
neither grouping partitions a pair. And the two ratio statistics the two-way form
is not written for use a one-way cluster bootstrap over cell lines, flagged where
they appear.

\paragraph{Biased evaluation} A version of Simpson's paradox. Because drugs act
at vastly different concentrations, most of the variance lies \emph{between}
drugs, so a model that learns nothing but each drug's mean explains most of it.
Their illustration is stark, a global Pearson of $0.90$ alongside an average
within-drug Pearson of $0.01$. \cref{sec:methods-metrics} explains why a
discrimination metric scored \emph{within} a comparison set is immune by
construction.

\paragraph{Splits} DrEval names four by application: Leave-Pairs-Out for
imputation, Leave-Cell-Lines-Out for personalised medicine, Leave-Drugs-Out for
drug design, Leave-Tissue-Out for
repurposing.\seealso{\cref{sec:res-generalisation} names the estimand this thesis
declines, and why}

Our holdout lands on one of those and deliberately misses another.
\texttt{unseen\_drug} is Leave-Drugs-Out. \texttt{unseen\_combo} is \emph{not}
Leave-Pairs-Out, and that name is reserved here for the estimand this thesis
declines: holding out a (drug, cell line) pair leaves the pair's own treated well
in training through its other $\approx 48$ cell lines, measuring within-atlas
matrix completion and not generalisation. Our split withholds the whole treatment
well, strictly the stronger of the two, and answers a narrower question, an
unseen well of a seen drug.

Citing Partin et al.~\cite{partin}, DrEval reports that good Leave-Pairs-Out and
Leave-Cell-Lines-Out performance is reachable by predicting each drug's
average response across the training cell lines. A strong result on those splits
therefore demands the per-drug average be measured explicitly, which is what
\cref{sec:res-lookup} discharges.

\paragraph{Missing ablations} Their remedy is randomisation and robustness
testing. The scramble ablation of \cref{sec:methods-controls} is one, specialised
to the drug. Hold every byte of the prompt fixed except the drug and see whether
anything downstream notices.

% ===========================================================================
\section{Cross-context transfer of drug responses}
\label{sec:lit-cmap}
\framedecl{-}

\begin{question}
An entire drug-discovery methodology rests on the premise that a drug's
characteristic transcriptional response travels between cellular contexts. The
premise is normally assumed. Here it is measured, with the attenuation from
sampling noise corrected.
\end{question}

\noindent
The Connectivity Map programme~\cite{cmap,l1000} built that methodology: a drug's
transcriptional response is taken to be stable enough across contexts that a
reference database can be queried by similarity to it. Its L1000
platform~\cite{l1000} ran the idea at scale, and its 978-gene landmark set
defines the gene panel used here, intersected with Tahoe's measured genes to give
$946$.

How much of that holds determines whether a conditional model has anything left
to predict once a per-drug constant is accounted for.\gloss{transfer
coefficient}{How much of a drug's residual in one context is recovered in
another, on a disattenuated similarity scale.}

\cref{sec:res-transfer} measures it on condition-level pseudobulk residuals. On a
disattenuated cross-context similarity scale the shared part is
\ci{\Tcoef = 0.553}{0.514}{0.593}, a little over half, far above the
structure-matched cross-line negative control ($-0.008$) and far below
reproduction. The earlier same-plate, different-drug control ($-0.017$) is not a
second matched control. It differs structurally from the estimand and was
voided.\corrmark{the same-plate control, voided in
\cref{sec:res-transfer}}

\paragraph{What the remainder is not} The remaining \ci{44.7\%}{40.7\%}{48.6\%}
is not thereby shown to be biology. The compared pairs differ in cell line, in
dose and in treatment well at once, so the figure is not a variance share and not
a drug$\times$cell-line interaction, and the section says so where it reports it.
What it delivers is a calibration of the transfer premise, not a verdict on it,
and we are aware of no prior estimate of that quantity for chemical single-cell
perturbation data.

% ===========================================================================
\section{Probing versus intervention}
\label{sec:lit-interp}
\framedecl{-}

\begin{question}
A probe can read a drug's identity out of a model's internals. Does that mean the
model uses it? To learn whether a property is used, change it and watch what the
model emits.
\end{question}

\noindent
Linear probing establishes that information is \emph{present} in a
representation, not that it is \emph{used}. Elazar et al.~\cite{gurnee} make the
case for keeping the two apart. Their amnesic-probing design is adopted here with
a matched control added, since an intervention that fails to move the prediction
proves nothing unless a comparable intervention does.

\paragraph{A sharper form of the same distinction} Gurnee et al.~\cite{workspace}
report that language models maintain a small privileged set of representations,
the \emph{J-space}, found by a Jacobian Lens, and argue it functions as a global
workspace. Roughly ten to twenty-five concepts are active at any moment,
typically under $10\%$ of activation variance. The defining property is causal.
Swapping one active
workspace vector for another changes the answer to match, while content outside
may be richly represented and still inert.

\paragraph{What this thesis takes} The design is the swap.
\cref{sec:res-mechanism} builds a domain analogue, substituting a drug's learned
representation for another drug's and measuring whether the prediction moves
toward the substituted drug. A property can be recoverable from the residual
stream and still sit outside whatever the readout consults, in which case making
it more explicit need not change what the model emits, an expectation
\cref{sec:res-epsilon} tests and \cref{sec:res-encoding} bounds.
The thesis also borrows the lens, in \cref{sec:res-jspace}.

% PLACEMENT ONLY. The box split 6 lines / 2 lines across the page break, and the
% 2-line tail arrived on the next page without the SCOPE LIMIT label, so it read
% as an unattributed indented paragraph. \budgettick reserves 5\baselineskip,
% which is not enough for this one: measured in the PDF the box draws 138.0bp,
% 9.3\baselineskip including its two rules and padding. 10 is the whole box.
% thesis_v5: the box's text was rewritten to point to sec:res-jspace and grew
% by a line or two, so the reservation grows with it.
\needspace{12\baselineskip}%
\begin{scopelimit}[carried into \cref{sec:res-mechanism} and \cref{sec:res-jspace}]
What the thesis does not do is test the workspace claim itself.
\cref{sec:res-jspace} measures how much of the change a drug swap makes lies
along the lens's token directions, which is geometric co-location and not use,
and no workspace vector is ever swapped. The substitution result
of \cref{sec:res-mechanism} is consistent with the workspace reading and
equally consistent with a readout that reads the drug at low gain for some
other reason; the co-location measurement is consistent with the first and
does not exclude the second. The framework is cited for the instruments it
justifies and the distinction it sharpens, not as a hypothesis this thesis
adjudicates.
\end{scopelimit}

\paragraph{Why allocate representations this way} Elhage et
al.~\cite{superposition} show that networks reserve dedicated directions for
features needed on every forward pass and push the rest into superposition, still
linearly recoverable but without an axis. A property required to emit any
plausible prediction should sit in the causally active set; one needed only for
occasional composition need not. That asymmetry motivates building the
instrument. It is not something the instrument returns.

% ===========================================================================
% NOTE TO SELF, removed from the printed title 2026-09-10: (REMEMBER TO CITE MERT KAAN THESIS)
\section{Optimal transport as a target constructor}
\label{sec:lit-ot}
\framedecl{-}


\begin{question}
Sequencing destroys the cell it reads, so no cell is seen both before and after
treatment and there is no paired training data. Where, then, does a training
target come from? The instrument is optimal transport, used not as the predictor
but as the thing that builds the target.
\end{question}

\noindent
There is no ground-truth correspondence between the control and treated
populations, only two clouds of points from the same well at different times, so
the correspondence must be \emph{inferred}.

\paragraph{Couplings} A deterministic Monge map, one destination per source
point, cannot split mass and is ill-posed for that reason. Kantorovich's
relaxation replaces it with a \emph{coupling}, a joint distribution
$\pi \in \Pi(\mu,\nu)$ with marginals $\mu$ and $\nu$,
\begin{equation}
  \min_{\pi \in \Pi(\mu,\nu)} \int_{\mathcal{X}\times\mathcal{Y}}
  c(x,y)\,\mathrm{d}\pi(x,y),
\end{equation}
for a ground cost $c$, typically squared Euclidean distance. A coupling may split
a source point across destinations, which makes the problem solvable for noisy,
overlapping empirical distributions. When $c$ is the $p$-th power of a metric the
optimal value defines $W_p(\mu,\nu)$, and $p=2$ is used here.

\paragraph{Regularisation}
Solved exactly, the discrete problem costs roughly $O(n^3 \log n)$, prohibitive
at single-cell scale. Adding a negative entropy term,
\begin{equation}
  \min_{\pi \in \Pi(\mu,\nu)} \; \langle C, \pi \rangle \;-\;
  \varepsilon H(\pi),
\end{equation}
renders it strictly convex, so the optimum is unique and takes a scaling form
recovered by the Sinkhorn--Knopp algorithm through alternating matrix--vector
products~\cite{cuturi}.

The parameter $\varepsilon$ is not merely a computational convenience. It
controls how sharply mass is assigned. As $\varepsilon \rightarrow 0$ the plan
concentrates toward a sparse coupling; without further conditions, that limit
need not send each control cell to one observed treated cell. As $\varepsilon$
grows the plan blurs, until every control cell couples equally to every treated
cell.

The \emph{barycentric projection} maps each control cell $x_i$ to the
$\pi$-weighted average of its partners,
\begin{equation}
  \hat{y}(x_i) \;=\; \frac{\sum_j \pi_{ij}\, y_j}{\sum_j \pi_{ij}},
\end{equation}
a convex combination inside the observed treated convex hull, not necessarily a
point on the data manifold. Small $\varepsilon$ therefore yields a target
dominated by a sparse set of observed treated cells, large $\varepsilon$ the
group mean, and a sweep interpolates continuously between two otherwise unrelated
target constructions.\seealso{\cref{sec:res-epsilon} sweeps this parameter and
scores the targets it produces}

\paragraph{A third use} Applied to biology this machinery mostly does
trajectory inference, as in Wad\-ding\-ton-OT~\cite{schiebinger}, or prediction,
taking the fitted map \emph{as} the predictor. The use made here is neither.
Transport supplies the correspondence from which a training target is built, and
what then has to generalise is the language model trained on that target, not the
coupling. Which is why what matters below is $\varepsilon$, the parameter
deciding which target the coupling defines.

\paragraph{Two practical points} The ground cost needs a space where Euclidean
distance is meaningful, and raw counts are not one, being sparse, overdispersed
and dominated by library size. Couplings are therefore computed in the released Tahoe scVI latent~\cite{scvi}. And because a full coupling costs quadratically, transport is solved per condition, coupling its treated cells to the controls shared within its cell-line/plate group. Both choices, the shipped scVI latent as the transport space and one coupling per condition against that condition's own controls, follow the optimal-transport pipeline of Kaan~\cite{mertkaan}, a companion thesis on the same atlas from which this construction was taken. A control cell from one cell line is not
a plausible counterfactual partner for a treated cell from another.
